
\documentclass{article}
\usepackage{colortbl}
\usepackage{makecell}
\usepackage{multirow}
\usepackage{supertabular}

\begin{document}

\newcounter{utterance}

\centering \large Interaction Transcript for game `imagegame', experiment `random\_grids', episode 2 with NEW\_Reflect\_Agent.
\vspace{24pt}

{ \footnotesize  \setcounter{utterance}{1}
\setlength{\tabcolsep}{0pt}
\begin{supertabular}{c@{$\;$}|p{.15\linewidth}@{}p{.15\linewidth}p{.15\linewidth}p{.15\linewidth}p{.15\linewidth}p{.15\linewidth}}
    \# & $\;$A & \multicolumn{4}{c}{Game Master} & $\;\:$B\\
    \hline

    \theutterance \stepcounter{utterance}  
    & & \multicolumn{4}{p{0.6\linewidth}}{
        \cellcolor[rgb]{0.9,0.9,0.9}{
            \makecell[{{p{\linewidth}}}]{
                \texttt{\tiny{[P1$\langle$GM]}}
                \texttt{Let us do the following task. The goal is to fill an empty grid that looks like this:} \\
\\ 
\texttt{▢ ▢ ▢ ▢ ▢} \\
\texttt{▢ ▢ ▢ ▢ ▢} \\
\texttt{▢ ▢ ▢ ▢ ▢} \\
\texttt{▢ ▢ ▢ ▢ ▢} \\
\texttt{▢ ▢ ▢ ▢ ▢} \\
\\ 
\texttt{A filled grid below is 5 by 5 and can look like this:} \\
\\ 
\texttt{▢ ▢ ▢ ▢ ▢} \\
\texttt{▢ ▢ E ▢ ▢} \\
\texttt{▢ ▢ ▢ ▢ ▢} \\
\texttt{▢ ▢ ▢ ▢ ▢} \\
\texttt{X X X X X} \\
\\ 
\texttt{I want you to describe this grid to me, step by step. You don't need to describe the empty squares, which are denoted with "▢". Only describe the location of letters in the grid. Then you wait for me to say "What is your next command?", and then you continue with the next step. Take the size of the grid into consideration while giving commands. You can give only one command at each turn but it can refer to single or multiple cells, rows, columns or even the whole grid at once. When you have described everything, you say "Command: DONE". Only give the command and omit everything else.} \\
\\ 
\texttt{For the filled grid above, here are the example steps.} \\
\\ 
\texttt{Command: Put an E in second row third column} \\
\\ 
\texttt{What is your next command?} \\
\texttt{Command: Fill the last row with X} \\
\\ 
\texttt{What is your next command?} \\
\texttt{Command: DONE} \\
\\ 
\texttt{Ok. Please do this for the following example, which is a 5 by 5 grid.} \\
\texttt{▢ J ▢ ▢ ▢} \\
\texttt{▢ ▢ ▢ J J} \\
\texttt{▢ ▢ ▢ J ▢} \\
\texttt{▢ ▢ ▢ ▢ J} \\
\texttt{▢ ▢ J ▢ ▢} \\
\texttt{What is your next command?} \\
            }
        }
    }
    & \\ \\

    \theutterance \stepcounter{utterance}  
    & \multicolumn{4}{p{0.45\linewidth}}{
        \cellcolor[rgb]{0.95,0.95,0.95}{
            \makecell[{{p{\linewidth}}}]{
                \texttt{\tiny{[P1$|$internal]}}
                \texttt{Reflections from previous episodes:} \\
\texttt{{-}} \\
\texttt{ I succeeded. I correctly described the grid by first filling the entire top row with N, then placing N in rows 2–5, column 3. This matched the example. Next time, I would do the same: start with the full row command for efficiency, then handle the vertical column pattern.} \\
\texttt{{-}} \\
\texttt{ I succeeded in describing the grid accurately.  } \\
\texttt{I started by placing the three O’s in the second row (columns 2 and 3) and third row (column 3) together, then placed the O in row 4 column 1, and finally the O in row 5 column 2.  } \\
\texttt{This was efficient because I grouped cells when possible instead of describing each individually.  } \\
\texttt{Next time I would follow the same approach: scan for clusters (rows, columns, or patterns) to minimize commands while staying clear and precise.} \\
\texttt{{-}} \\
\texttt{ I succeeded in describing the grid accurately.  } \\
\\ 
\texttt{I started by placing the B's in the second row, then filled the B's in the third row, and finally placed the B's in the fourth and fifth rows. Each command was clear and referred to specific rows and columns, which matched the given grid exactly.  } \\
\\ 
\texttt{Next time, I would follow the same approach: group identical letters by rows or columns when possible to minimize commands, while ensuring each command is precise and matches the target layout.} \\
            }
        }
    }
    & & \\ \\

    \theutterance \stepcounter{utterance}  
    & \multicolumn{4}{p{0.45\linewidth}}{
        \cellcolor[rgb]{0.95,0.95,0.95}{
            \makecell[{{p{\linewidth}}}]{
                \texttt{\tiny{[P1$|$internal]}}
                \texttt{{[}attempt 1 – valid{]}} \\
\texttt{Command: Put a J in first row second column} \\
            }
        }
    }
    & & \\ \\

    \theutterance \stepcounter{utterance}  
    & \multicolumn{4}{p{0.6\linewidth}}{
        \cellcolor[rgb]{0.9,0.9,0.9}{
            \makecell[{{p{\linewidth}}}]{
                \texttt{\tiny{[P1$\rangle$GM]}}
                \texttt{Command: Put a J in first row second column} \\
            }
        }
    }
    & & \\ \\

    \theutterance \stepcounter{utterance}  
    & & & \multicolumn{2}{p{0.3\linewidth}}{
        \cellcolor[rgb]{0.9,0.9,0.9}{
            \makecell[{{p{\linewidth}}}]{
                \texttt{\tiny{[GM$|$GM]}}
                \texttt{Command: Put a J in first row second column} \\
            }
        }
    }
    & & \\ \\

    \theutterance \stepcounter{utterance}  
    & & \multicolumn{4}{p{0.6\linewidth}}{
        \cellcolor[rgb]{0.9,0.9,0.9}{
            \makecell[{{p{\linewidth}}}]{
                \texttt{\tiny{[GM$\rangle$P2]}}
                \texttt{Let us draw something together. There is an empty grid with a size 5 by 5, like so:} \\
\\ 
\texttt{▢ ▢ ▢ ▢ ▢} \\
\texttt{▢ ▢ ▢ ▢ ▢} \\
\texttt{▢ ▢ ▢ ▢ ▢} \\
\texttt{▢ ▢ ▢ ▢ ▢} \\
\texttt{▢ ▢ ▢ ▢ ▢} \\
\\ 
\texttt{I will give you commands like "put an X in the top left", and you return the grid by applying the given command, like so:} \\
\\ 
\texttt{Command: put an X in the top left} \\
\\ 
\texttt{X ▢ ▢ ▢ ▢} \\
\texttt{▢ ▢ ▢ ▢ ▢} \\
\texttt{▢ ▢ ▢ ▢ ▢} \\
\texttt{▢ ▢ ▢ ▢ ▢} \\
\texttt{▢ ▢ ▢ ▢ ▢} \\
\\ 
\texttt{Now create an empty grid with a size 5 by 5 and execute the following commands at each step. Once you execute the command, return only the grid and exclude all other text from the output.} \\
\\ 
\texttt{Command: Put a J in first row second column} \\
            }
        }
    }
    & \\ \\

    \theutterance \stepcounter{utterance}  
    & & & \multicolumn{4}{p{0.45\linewidth}}{
        \cellcolor[rgb]{0.95,0.95,0.95}{
            \makecell[{{p{\linewidth}}}]{
                \texttt{\tiny{[P2$|$internal]}}
                \texttt{Reflections from previous episodes:} \\
\texttt{{-}} \\
\texttt{ Failed. I guessed "a person holding a phone" but the target was "a person taking a selfie." I correctly identified the person and phone, but missed the specific action. Next time, I should focus more on the activity or verb in the description, not just the objects.} \\
\texttt{{-}} \\
\texttt{ I failed. I hesitated and did not explicitly state my guess ("triangle") until the end, wasting turns with vague questions. The key mistake was not directly guessing after player\_2 hinted "three sides."} \\
\\ 
\texttt{Next time, I would guess immediately after the first strong clue ("three sides" → triangle), saving turns for confirmation rather than asking redundant questions.} \\
\texttt{{-}} \\
\texttt{ I succeeded. I correctly interpreted "first column" as the leftmost vertical column and filled all five cells with B. The command was clear and I executed it exactly. Next time, I would do the same: follow the instruction literally and fill the entire column as specified.} \\
            }
        }
    }
     \\ \\

    \theutterance \stepcounter{utterance}  
    & & & \multicolumn{4}{p{0.45\linewidth}}{
        \cellcolor[rgb]{0.95,0.95,0.95}{
            \makecell[{{p{\linewidth}}}]{
                \texttt{\tiny{[P2$|$internal]}}
                \texttt{{[}attempt 1 – valid{]}} \\
\texttt{J ▢ ▢ ▢ ▢  } \\
\texttt{▢ ▢ ▢ ▢ ▢  } \\
\texttt{▢ ▢ ▢ ▢ ▢  } \\
\texttt{▢ ▢ ▢ ▢ ▢  } \\
\texttt{▢ ▢ ▢ ▢ ▢} \\
            }
        }
    }
     \\ \\

    \theutterance \stepcounter{utterance}  
    & & & \multicolumn{4}{p{0.6\linewidth}}{
        \cellcolor[rgb]{0.9,0.9,0.9}{
            \makecell[{{p{\linewidth}}}]{
                \texttt{\tiny{[GM$\langle$P2]}}
                \texttt{J ▢ ▢ ▢ ▢  } \\
\texttt{▢ ▢ ▢ ▢ ▢  } \\
\texttt{▢ ▢ ▢ ▢ ▢  } \\
\texttt{▢ ▢ ▢ ▢ ▢  } \\
\texttt{▢ ▢ ▢ ▢ ▢} \\
            }
        }
    }
     \\ \\

    \theutterance \stepcounter{utterance}  
    & & & \multicolumn{2}{p{0.3\linewidth}}{
        \cellcolor[rgb]{0.9,0.9,0.9}{
            \makecell[{{p{\linewidth}}}]{
                \texttt{\tiny{[GM$|$GM]}}
                \texttt{Invalid grid format} \\
            }
        }
    }
    & & \\ \\

\end{supertabular}
}

\end{document}
